% !TeX root = ../navier-stokes-manuscript.tex
% ============================================================
% 1.2 Where the Problem Is Today
% ============================================================
\subsection{The Classical Estimates Before Closure}\label{sub:HistoricalEstimates}

Local theory gives the smooth fluid a beginning, so the first serious question is what the equations preserve after that beginning.
The first answer is kinetic energy, and almost every later regularity result can be read as an attempt to recover the derivative control that energy leaves behind.

\divider{Energy and the First Uncontrolled Derivative}

For a smooth solution, multiply the momentum equation by \(u\) and integrate over space.
Incompressibility removes nonlinear transport from the total energy because
\[
  \int_{\mathbb R^3}(u\cdot\nabla)u\cdot u\,dx
  =
  \frac12\int_{\mathbb R^3}u\cdot\nabla|u|^2\,dx
  =0.
\]
It removes the pressure contribution for the same divergence-free reason:
\[
  \int_{\mathbb R^3}\nabla p\cdot u\,dx
  =
  -\int_{\mathbb R^3}p\,\nabla\cdot u\,dx
  =0.
\]
Viscosity is the surviving term, and integration by parts gives
\[
  \nu\int_{\mathbb R^3}\Delta u\cdot u\,dx
  =
  -\nu\lVert\nabla u\rVert_{L^2}^2.
\]
The resulting identity is
\[
  \lVert u(t)\rVert_{L^2}^2
  +
  2\nu\int_0^t\lVert\nabla u(s)\rVert_{L^2}^2\,ds
  =
  \lVert u_0\rVert_{L^2}^2.
\]
The velocity has an \(L^2\) ceiling at every time, and viscosity pays for the squared \(L^2\) gradient after it is accumulated in time.

Leray found a way to keep this estimate when smoothness is no longer assumed.
Regularization, compactness, and the energy inequality produce a global weak solution satisfying \parencite{Leray1934}
\[
  u\in
  L^\infty(0,T;L^2(\mathbb R^3))
  \cap
  L^2(0,T;H^1(\mathbb R^3))
\]
on every finite interval.
This is genuine global existence at the energy level; the energy inequality alone does not settle uniqueness or spontaneous regularity for arbitrary three-dimensional Leray solutions.

The Sobolev embedding
\[
  H^1(\mathbb R^3)\hookrightarrow L^6(\mathbb R^3)
\]
turns Leray's spatial derivative into
\[
  u\in L^2(0,T;L^6(\mathbb R^3)).
\]
That extra integrability is immediately spent on the nonlinear transport.
For example,
\[
  \lVert(u\cdot\nabla)u\rVert_{L^{3/2}}
  \leq
  \lVert u\rVert_{L^6}
  \lVert\nabla u\rVert_{L^2},
\]
so the transport term is integrable in time with values in \(L^{3/2}\).
This is strong enough to keep the equation meaningful in the weak construction.
It still supplies no time-uniform ceiling on \(\lVert\nabla u(t)\rVert_{L^2}\), because a sequence of higher and shorter peaks can have a finite squared time integral.

Vorticity shows why that opening is specifically dangerous in three dimensions.
Set
\[
  \omega:=\nabla\times u.
\]
For the divergence-free decaying fields considered here,
\[
  \lVert\nabla u\rVert_{L^2}^2
  =
  \lVert\omega\rVert_{L^2}^2,
\]
so the dissipated gradient is also the enstrophy.
Taking the curl of the momentum equation gives
\[
  \partial_t\omega+(u\cdot\nabla)\omega
  =
  (\omega\cdot\nabla)u+\nu\Delta\omega.
\]
The new term \((\omega\cdot\nabla)u\) stretches and tilts vorticity with the velocity gradient generated by the same flow.
Its energy balance is
\[
  \frac12\frac{d}{dt}\lVert\omega\rVert_{L^2}^2
  +
  \nu\lVert\nabla\omega\rVert_{L^2}^2
  =
  \int_{\mathbb R^3}(\omega\cdot\nabla)u\cdot\omega\,dx.
\]
The stretching integral has no fixed sign, and the energy identity does not force it to remain below viscous dissipation at each time.
The first estimate reaches the whole time interval at a level below the derivative control needed to restart the smooth solution.

\divider{What Is Strong Enough to Continue the Solution}

The natural way to ask for several derivative levels at once is a Sobolev norm.
For an integer \(m\geq0\), write
\[
  \lVert u\rVert_{H^m}^2
  :=
  \sum_{|\alpha|\leq m}
  \lVert\partial^\alpha u\rVert_{L^2}^2,
\]
with the usual Fourier definition when the order is fractional.
In three dimensions, Sobolev embedding gives \parencite[Theorem~7.28]{Cerda2010}
\[
  H^s(\mathbb R^3)\hookrightarrow C^k(\mathbb R^3)
  \qquad
  \text{when}
  \qquad
  s>k+\frac32.
\]
In particular,
\[
  s>\frac52
  \qquad\Longrightarrow\qquad
  H^s(\mathbb R^3)\hookrightarrow W^{1,\infty}(\mathbb R^3).
\]
This is the first level at which one Sobolev norm directly controls the pointwise velocity gradient that multiplies every higher differentiated energy estimate.

Suppose a smooth solution on \([0,T)\) satisfies
\[
  \sup_{0\leq t<T}\lVert u(t)\rVert_{H^s}
  \leq M
  \qquad
  \text{for some }s>\frac52.
\]
For each fixed finite \(m\geq s\), differentiate the equation spatially through order \(m\), pair with the corresponding derivative of \(u\), and sum the resulting identities.
The product rule and incompressibility give the standard estimate
\[
  \frac12\frac{d}{dt}\lVert u\rVert_{H^m}^2
  +
  \nu\lVert\nabla u\rVert_{H^m}^2
  \leq
  C_m\lVert\nabla u\rVert_{L^\infty}
  \lVert u\rVert_{H^m}^2.
\]
The assumed \(H^s\) ceiling controls the coefficient on the right, and Gronwall's inequality propagates the chosen \(H^m\) norm from the smooth datum.
The constant may depend on \(m\), because smoothness asks for every finite order in turn rather than one constant for the entire infinite collection.

Pressure does not add an independent unknown at the next stage.
Taking the divergence of the momentum equation gives
\[
  -\Delta p
  =
  \sum_{i,j=1}^3\partial_i\partial_j(u_i u_j),
\]
with decay fixing the pressure normalization.
Once the spatial derivatives of \(u\) are controlled, this elliptic relation recovers the corresponding spatial derivatives of \(p\), and
\[
  \partial_tu
  =
  \nu\Delta u-(u\cdot\nabla)u-\nabla p
\]
then recovers the first time derivative.
Repeating the same two operations recovers each higher mixed spatial and temporal derivative.
One uniform Sobolev bound above \(5/2\) has enough leverage to return the complete smooth velocity--pressure pair.

It also gives the missing time needed to cross an endpoint.
Local well-posedness in \(H^s\) supplies a positive lifespan determined by \(s\), \(\nu\), and the size of the initial \(H^s\) norm \parencite{FujitaKato1964,Kato1984}.
Choose \(t_0<T\), regard \(u(\cdot,t_0)\) as a fresh initial slice, and restart the local theory from that slice.
If \(\lVert u(\cdot,t_0)\rVert_{H^s}\leq M\) uniformly, the restarted lifespans have a common positive lower bound, so a sufficiently late \(t_0\) carries the same solution past \(T\).
Uniqueness on the overlap joins the restarted piece to the original motion.
Hence a finite maximal time forces
\[
  \limsup_{t\uparrow T_*}
  \lVert u(t)\rVert_{H^s}
  =\infty
  \qquad
  \text{for every }s>\frac52.
\]
The distance from Leray's estimate to continuation is now exact: time-integrated \(H^1\) control must somehow become time-uniform control at a Sobolev level above \(5/2\).

\divider{Critical Scale and Singular Concentration}

The equations themselves show where such a conversion can lose information.
If \((u,p)\) is a solution, then
\[
  u_\lambda(x,t)
  :=
  \lambda u(\lambda x,\lambda^2t),
  \qquad
  p_\lambda(x,t)
  :=
  \lambda^2p(\lambda x,\lambda^2t)
\]
has the same equation at a different spatial and temporal scale.
A quantity left unchanged by this transformation sees an intensifying structure with the same strength no matter how small that structure becomes.

The Ladyzhenskaya--Prodi--Serrin criteria locate exactly this balance between spatial intensity and temporal persistence \parencite{Prodi1959,Serrin1962,Serrin1963,Ladyzhenskaya1967}.
They give regularity through \(T\) when
\[
  u\in L^p(0,T;L^q(\mathbb R^3)),
  \qquad
  \frac{2}{p}+\frac{3}{q}\leq1,
  \qquad
  q>3.
\]
The spatial exponent measures how tightly velocity may concentrate, and the time exponent measures how long that concentration may persist.
Equality in the exponent relation is unchanged by Navier--Stokes scaling.

The spatial endpoint of this family is \(L^\infty_tL^3_x\), because
\[
  \lVert u_\lambda(t)\rVert_{L^3}
  =
  \lVert u(\lambda^2t)\rVert_{L^3}.
\]
Escauriaza, Seregin, and \v{S}ver{\'a}k proved that
\[
  u\in L^\infty(0,T;L^3(\mathbb R^3))
\]
still carries a solution regularly through \(T\) \parencite{EscauriazaSereginSverak2003}.
Their rescaling argument turns a first singularity into an ancient limiting solution and uses backward uniqueness to rule that limit out.
Thus a finite endpoint cannot hide behind a bounded critical \(L^3\) norm.

The local theory of singular concentration reaches a complementary conclusion.
For \(z_0=(x_0,t_0)\), write the parabolic cylinder
\[
  Q_r(z_0)
  :=
  B_r(x_0)\times(t_0-r^2,t_0).
\]
Caffarelli, Kohn, and Nirenberg proved that sufficiently small scale-invariant velocity and pressure quantities on \(Q_r(z_0)\) force regularity on a smaller cylinder \parencite{CaffarelliKohnNirenberg1982}.
A singular point must retain a definite scale-invariant concentration in every sufficiently small parabolic neighbourhood.
Their covering argument confines the possible singular set to zero one-dimensional parabolic Hausdorff measure, but a single point would still end global smoothness.

The endpoint result also has a direct blow-up consequence.
For the smooth compactly supported whole-space data in Seregin's theorem, a finite maximal time satisfies \parencite{Seregin2012}
\[
  \lim_{t\uparrow T_*}
  \lVert u(t)\rVert_{L^3(\mathbb R^3)}
  =\infty.
\]
Tao's quantitative refinement forces at least a positive power of a triple logarithm of the inverse time remaining along infinitely many times approaching the endpoint \parencite{Tao2021Quantitative}.
That forced rate is extremely slow compared with any power law, so much faster concentration remains compatible with all known critical lower bounds.

\divider{What the Classical Estimates Leave Open}

Taken on their own, the classical estimates put any finite maximal time in a narrow but still inhabited gap.
Kinetic energy stays bounded, every continuation-grade Sobolev norm loses finite control, the critical \(L^3\) norm diverges, and every singular point retains scale-invariant concentration at arbitrarily small parabolic scales.
None of those statements prevents the concentration itself.

The reason can be seen directly from scaling:
\[
  \lVert u_\lambda(t)\rVert_{L^3}
  =
  \lVert u(\lambda^2t)\rVert_{L^3},
  \qquad
  \lVert u_\lambda(t)\rVert_{L^2}^2
  =
  \lambda^{-1}
  \lVert u(\lambda^2t)\rVert_{L^2}^2.
\]
Moving a structure to a shorter length scale preserves its critical \(L^3\) reading while reducing the \(L^2\) energy carried by the rescaled piece.
The global energy ceiling cannot become a scale-independent ceiling on that critical concentration.

Time also contracts with scale.
If a schematic cascade halves its active length at each stage, its parabolic durations are proportional to
\[
  1,\ \frac14,\ \frac1{16},\ \ldots,
\]
and these durations fit inside a finite total time.
This is the Zeno-cascade picture: successively smaller and faster activity is not excluded by the known estimates simply because each stage consumes less time.
It describes room left open by the estimates and does not construct a singular Navier--Stokes fluid.

The classical estimates identify both the continuation-grade norm and the scale-critical behavior forced by a finite endpoint, but they do not supply terminal-uniform control of that norm.
